\documentclass[12pt,a4paper]{article}
\usepackage[utf8]{inputenc}
\usepackage[german]{babel}
\usepackage{amsmath}
\usepackage{amsfonts}
\usepackage{amssymb}
\usepackage{graphicx}
\usepackage{hyperref}
\usepackage{booktabs}
\usepackage{siunitx}
\usepackage{verbatim}

\title{Bachelorarbeit: Low-Power IoT-Skala}
\author{Thomas Rocheteau}
\date{Februar 2026}

\begin{document}

	\maketitle

	\tableofcontents
	\newpage

	% ============================================================
	\section{Einleitung}
	% ============================================================

	\subsection{Ausgangslage}
	% Hier Text einfügen

	\subsection{Zielsetzung}
	% Hier Text einfügen

	\subsection{Meilensteine}

	\begin{enumerate}
		\item \textbf{Hardware-Entwicklung:} Custom ESP32 Low-Power Board vs. kommerzielle Lösungen.
		\item \textbf{Quality of Life Improvements:} Implementierung von Bedienelementen (Taster) sowie Software-Funktionen (Tara, Kalibrierung).
		\item \textbf{Kommunikationsprotokolle:} Evaluation von ESP-NOW vs. BLE.
	\end{enumerate}

	% ============================================================
	\section{Requirements Engineering}
	% ============================================================
	% Hier Text einfügen

	% ============================================================
	\section{Technische Ausgangslage}
	% ============================================================

	\subsection{Design Space Exploration}

	\subsubsection{MCU-Auswahl: ESP32-S3 vs. ESP32-C3}
	\begin{itemize}
		\item \textbf{ESP32-S3:} Verfügt über einen Ultra-Low-Power (ULP) Co-Prozessor. Ermöglicht präzises Abtasten und Detektion von Gewichtsänderungen im Tiefschlaf . Datasheet sagt 8uA\cite{rnt_supermini}\cite{s3_datasheet}.
		\item \textbf{ESP32-C3:} Intervall-basiertes Abtasten (z. B. alle 5 Minuten) durch periodisches Aufwachen des Hauptprozessors. Datasheet sagt 5uA \cite{c3_datasheet}
		- ESP32-C3 integrates two 12-bit SAR ADCs.
	\end{itemize}

	\subsubsection{Evaluation kommerzieller Boards}
	\begin{itemize}
		\item \textbf{ESP32-C3 Super Mini:} Spezifiziert mit einem Low-Power-Stromverbrauch von ca. $43\,\mu A$ (im Deep Sleep).
	\end{itemize}

	\subsubsection{Kommunikationsstrategie}
	Untersuchung der Vor- und Nachteile von \textbf{ESP-NOW} gegenüber \textbf{BLE} (Bluetooth Low Energy).

	% ============================================================
	\section{Implementation Hardware}
	% ============================================================

	\subsection{Dimensionierung der Spannungsversorgung (Buck-Boost)}
	Für die Systemspannung wurde der hocheffiziente Buck-Boost-Konverter \textit{TPS63802} gewählt. Dieser ermöglicht es, die schwankende Spannung der 18650-Lithium-Ionen-Zelle ($3,0\text{V}$ bis $4,2\text{V}$) stabil auf das Zielniveau zu wandeln.

	\subsubsection{Berechnung des Feedback-Netzwerks}
	Die Ausgangsspannung $V_{OUT}$ wird über einen externen Spannungsteiler am Feedback-Pin (FB) eingestellt. Die interne Referenzspannung beträgt beim TPS63802 konstant $V_{FB} = 500\text{mV}$. Die Berechnung erfolgt nach der Formel:

	\begin{equation}
		V_{OUT} = V_{FB} \cdot \left(1 + \frac{R_1}{R_2}\right)
	\end{equation} \cite{TPS63802}

	In der Testphase werden zwei Konfigurationen evaluiert, um das Optimum zwischen Energieeffizienz und Signalqualität der Wägezelle zu ermitteln:

	\begin{table}[h]
		\centering
		\caption{Widerstandswerte für unterschiedliche Ausgangsspannungen}

		\begin{tabular}{@{}llll@{}}
			\toprule
			Zielspannung $V_{OUT}$ & $R_1$ (E96-Reihe) & $R_2$ (Standard) & Berechnetes $V_{OUT}$ \\ \midrule
			$3,3\text{V}$ & $511\,\text{k}\Omega$ & $91\,\text{k}\Omega$ & $3,307\text{V}$ \\
			$3,0\text{V}$ & $500\,\text{k}\Omega$ & $100\,\text{k}\Omega$ & $3,000\text{V}$ \\ \bottomrule
		\end{tabular}
	\end{table}

	\subsubsection{Optimierung der Peripheriebeschaltung}
	Auf die Implementierung des \textit{Power-Good} (PG) Indikators des TPS63802 wurde bewusst verzichtet:

	\begin{itemize}
		\item \textbf{Reduktion der Komplexität:} Da der ESP32-C3 über interne Brown-out-Detektoren verfügt, ist eine externe Überwachung der Spannungsstabilität über den PG-Pin redundant.
		\item \textbf{Minimierung des Ruhestroms:} Der Verzicht auf den erforderlichen Pull-up-Widerstand am Open-Drain-Ausgang des PG-Pins eliminiert potenzielle Leckströme im Fehlerfall oder während der Startphase.
		\item \textbf{Platzbedarf:} Die Einsparung dieses Schaltungsteils erlaubt ein kompakteres Routing der sensiblen Feedback-Leitungen, was die EMV-Eigenschaften des Buck-Boost-Konverters verbessert.
	\end{itemize}

	\subsubsection{Versuchsaufbau: $3,3\text{V}$ vs. $3,0\text{V}$}
	Im Rahmen der Validierung wird untersucht, ob eine Absenkung der Systemspannung auf $3,0\text{V}$ signifikante Vorteile bringt. Folgende Hypothesen werden geprüft:

	\begin{enumerate}
		\item \textbf{Energieeinsparung:} Da die Leistungsaufnahme vieler Komponenten quadratisch mit der Spannung korreliert ($P = U^2 / R$), könnte $3,0\text{V}$ die Batterielaufzeit verlängern.
		\item \textbf{Stabilität des Mikrocontrollers:} Es wird verifiziert, ob der ESP32-C3 bei Stromspitzen während der WiFi-Übertragung unter $3,0\text{V}$ stabil arbeitet.
		\item \textbf{Signal-Rausch-Verhältnis (SNR):} Eine geringere Versorgungsspannung am HX711-Verstärker reduziert die Erregerspannung der Wägezelle. Es ist zu klären, ob dies die Messgenauigkeit der Waage negativ beeinflusst.
	\end{enumerate}

	Zur Minimierung des Ruhestroms wurden bewusst hochohmige Widerstände im Megaohm-Bereich gewählt, um den Querstrom des Teilers unter $5\mu\text{A}$ zu halten.

	\subsection{Energieversorgung (Power Supply)}
	Die Dimensionierung der Energieversorgung basiert auf folgendem Setup:
	\begin{itemize}
		\item Konfiguration: $3 \times 18650$ Akkumulatoren
		\item Kapazität: $3 \times 3000\,mAh = 9000\,mAh$
		\item Energiegehalt: $11,1\,Wh$
	\end{itemize}

	\subsection{Einfluss der Entladekurve auf die Effizienz}

	Ein kritischer Aspekt beim Vergleich zwischen Linearreglern (LDO) und Schaltreglern (Buck-Boost) ist das Verhalten über den gesamten Spannungsverlauf der Batterie. Während der Buck-Boost-Converter theoretisch eine nahezu konstant hohe Effizienz bietet, ist die Effizienz des LDO direkt an das Spannungsgefälle zwischen Eingangs- und Ausgangsspannung gekoppelt.

	\subsubsection{Die LDO-Effizienz-Kurve}
	Die Effizienz $\eta_{LDO}$ berechnet sich näherungsweise aus dem Verhältnis von Ausgangsspannung $U_{out}$ zu Eingangsspannung $U_{in}$:

	\begin{equation}
		\eta_{LDO} \approx \frac{U_{out}}{U_{in}}
	\end{equation}

	Daraus ergeben sich für eine Li-Ion-Zelle (Typ 18650) bei einer geregelten Ausgangsspannung von $3,3\,V$ folgende Werte:

	\begin{itemize}
		\item \textbf{Bei voller Batterie ($4,2\,V$):} $\eta \approx 78,5\,\%$
		\item \textbf{Bei Nominalspannung ($3,7\,V$):} $\eta \approx 89,1\,\%$
		\item \textbf{Kurz vor Entladeschluss ($3,4\,V$):} $\eta \approx 97,0\,\%$
	\end{itemize}

	\subsubsection{Dynamik des Break-Even-Punkts}
	Da die Effizienz des LDO mit sinkender Batteriespannung steigt, verschiebt sich auch der Break-Even-Punkt zuungunsten des Buck-Boost-Converters.

	\begin{enumerate}
		\item \textbf{Frühphase ($4,2\,V$):} Der Buck-Boost spart ca. $15\,mA$ im Aktiv-Modus. Der Break-Even-Punkt liegt bei ca. $54$ Minuten.
		\item \textbf{Mittlere Phase ($3,7\,V$):} Die Differenz $\Delta I_{active}$ sinkt auf ca. $8\,mA$. Der resultierende Break-Even-Punkt verschiebt sich wie folgt:
	\end{enumerate}

	\begin{equation}
		t_{sleep} = \frac{8.000\,\mu As \cdot 2\,s}{9\,\mu A} \approx 1.777\,s \approx \mathbf{29,6\,Minuten}
	\end{equation}

	\textbf{Interpretation:} Je weiter die Batterie entladen wird, desto kürzer müssten die Intervalle zwischen den Aktivphasen sein, damit sich der Einsatz eines Buck-Boost-Converters energetisch amortisiert. Da das System aufgrund der geringen Solarleistung von $19\,\mu A$ vorwiegend im mittleren oder unteren Spannungsbereich operieren wird, schwindet der theoretische Vorteil des Schaltreglers in der praktischen Anwendung nahezu vollständig.

	\subsubsection{Zusammenfassende Bewertung für die Hardwarewahl}
	Die Wahl fiel auf den \textbf{RT9080-33GJ5}, da dieser in der Gesamtbetrachtung drei wesentliche Vorteile für das spezifische Lastprofil vereint:

	\begin{itemize}
		\item \textbf{Minimaler statischer Verlust:} Mit einem Ruhestrom $I_q$ von nur $2\,\mu A$ wird die begrenzte Solar-Energieernte minimal belastet.
		\item \textbf{Dynamische Effizienz:} Im wahrscheinlichen Spannungsfenster ($3,4\,V$ bis $3,8\,V$) arbeitet der LDO mit einer Effizienz von über $85\,\%$, was unter Berücksichtigung der Komplexität vergleichbar mit einem Schaltregler bei geringen Lasten ist.
		\item \textbf{Optimiertes Dropout-Verhalten:} Die extrem niedrige Dropout-Spannung ermöglicht den stabilen Betrieb der $3,3\,V$-Elektronik bis zu einer Batteriespannung von ca. $3,35\,V$. Dies maximiert die nutzbare Kapazität der 18650-Zelle.
	\end{itemize}

	\subsection{Sensorik: Wägezellen und ADC}

	\subsubsection{Wägezellen (YZC-133)}

	Als Kraftsensoren kommen vier Wägezellen des Typs YZC-133 zum Einsatz \cite{yzc133_datasheet}.
	Jede Zelle ist für eine Nennlast von \SI{1}{\kilogram} ausgelegt und enthält intern eine vollständige Wheatstone-Brücke aus vier Dehnungsmessstreifen (DMS).
	Die relevanten elektrischen Kenngrössen sind in Tabelle~\ref{tab:yzc133} zusammengefasst.

	\begin{table}[h]
		\centering
		\caption{Elektrische Spezifikation YZC-133}
		\label{tab:yzc133}
		\begin{tabular}{@{}ll@{}}
			\toprule
			Parameter & Wert \\ \midrule
			Nennlast                & \SI{1}{\kilogram} \\
			Nennkennwert (Empfindlichkeit) & $1{,}0 \pm 0{,}15\,\text{mV/V}$ \\
			Nullpunktsignal         & $\pm 0{,}1\,\text{mV/V}$ \\
			Kriechfehler            & $0{,}03\,\%\,\text{F.S.} / \SI{30}{\minute}$ \\
			Eingangsimpedanz        & $1066 \pm 20\,\Omega$ \\
			Ausgangsimpedanz        & $1000 \pm 5\,\Omega$ \\ \bottomrule
		\end{tabular}
	\end{table}

	Bei \SI{3{,}3}{\volt} Erregerspannung liegen beide Ausgangspins der Brücke auf einem Gleichtaktpotenzial von ca.\ \SI{1{,}65}{\volt} ($V_{\text{ex}}/2$).
	Das Nutzsignal ist die \emph{Differenz} zwischen diesen beiden Pins: bei Vollast nominell $1{,}0\,\text{mV/V} \times 3{,}3\,\text{V} = \SI{3{,}3}{\milli\volt}$.
	Der differentielle Eingang des ADS1234 erfasst genau diese Differenz und unterdrückt das Gleichtaktsignal durch seine hohe Gleichtaktunterdrückung (CMRR).
	Die vier Zellen sind in X-Formation an den Ecken der Waagplattform montiert.

	\subsubsection{ADS1234 -- 4-Kanal-24-Bit-ADC}

	Für die Signalerfassung wird der ADS1234 von Texas Instruments eingesetzt \cite{ads1234_datasheet}.
	Der Chip ist ein 24-Bit-Sigma-Delta-ADC mit vier differenziellen Eingangskanälen und einem integrierten programmierbaren Verstärker (PGA) mit den Verstärkungsfaktoren 1, 2, 64 und 128.
	Mit PGA-Gain 128 ergibt sich für das \SI{3{,}3}{\milli\volt}-Signal der YZC-133 ein verstärktes Differenzsignal von ca. \SI{422}{\milli\volt}, das den Eingangsbereich des ADC gut ausnutzt.
	Die wählbare Ausgabedatenrate beträgt 10~SPS oder 80~SPS; im vorliegenden System wird 10~SPS verwendet, da dies das Eingangsrauschen durch die längere Integrationszeit maximally reduziert.

	\subsubsection{Von einem auf vier ADC-Kanäle}

	In der Vorarbeit zu dieser Bachelorarbeit waren alle vier Wägezellen in einer gemeinsamen Wheatstone-Brücke zusammengeschaltet und lieferten ihr summiertes Signal an einen einzigen ADC-Kanal.
	In dieser Arbeit wird jede Wägezelle einzeln an einen eigenen Eingangskanal des ADS1234 geführt.
	Dieser Wechsel bringt zwei wesentliche Vorteile:

	\paragraph{Ecklastkompensation}
	Mit einem einzigen Summensignal ist keine positionsabhängige Korrektur möglich: Das System kann nicht unterscheiden, ob eine Abweichung durch einen exzentrischen Lastangriffspunkt oder durch eine tatsächliche Massenänderung entsteht.
	Durch die individuellen Kanäle lässt sich jede Wägezelle separat kalibrieren und der dominante Sensor pro Lastposition gezielt korrigieren (siehe Abschnitt~\ref{sec:kalibrierung}).

	\paragraph{Rauschverhalten und individuelle Kalibrierung}
	Die Empfindlichkeit der YZC-133 streut mit $\pm 15\,\%$ zwischen den Exemplaren.
	In der Einzelkanal-Lösung mittelte sich diese Streuung über die parallele Zusammenschaltung heraus, konnte aber nicht pro Sensor korrigiert werden.
	Mit vier separaten Kanälen wird die Empfindlichkeitsdifferenz jeder Wägezelle durch den individuell bestimmten Skalierungsfaktor (\texttt{spanFactor}) in der Software exakt kompensiert.
	Zusätzlich verbessert die digitale Summation vier statistisch unkorrelierter ADC-Messungen das Signal-Rausch-Verhältnis gegenüber einer einzigen Messung um den Faktor $\sqrt{4} = 2$ ($+6\,\text{dB}$), da sich das Rauschen der Kanäle inkohärent addiert, während sich die Nutzsignale kohärent addieren.

	% ============================================================
	\section{Implementation Software}
	% ============================================================

	\subsection{ADS1234 -- Ansteuerung und Kalibrierung}
	\label{sec:kalibrierung}

	\subsubsection{Überblick}

	Der ADS1234 ist ein 24-Bit-Sigma-Delta-ADC von Texas Instruments, der speziell für den Einsatz mit Brückensensoren wie Dehnungsmessstreifen (DMS) ausgelegt ist \cite{ads1234_datasheet}.
	Er stellt vier unabhängige Eingangskanäle bereit, die über zwei digitale Adressleitungen (\texttt{A0}, \texttt{A1}) ausgewählt werden.
	Die Kommunikation mit dem Mikrocontroller erfolgt über ein synchrones serielles Protokoll mit den Leitungen \texttt{DRDY/DOUT} (kombiniertes Datenbereit- und Datensignal), \texttt{SCLK} (Takt) sowie \texttt{PDWN} (Power-Down/Reset).

	\begin{table}[h]
		\centering
		\caption{Pin-Belegung ADS1234 -- ESP32}
		\begin{tabular}{@{}lll@{}}
			\toprule
			Signal & GPIO & Funktion \\ \midrule
			\texttt{DRDY/DOUT} & 2  & Datenbereit-Anzeige und serieller Datenausgang \\
			\texttt{SCLK}      & 15 & Serieller Takt (Master-Ausgang) \\
			\texttt{PDWN}      & 26 & Power-Down / Reset \\
			\texttt{DMS\_PWR}  & 17 & Versorgung der Wägezellen (schaltbar) \\
			\texttt{A0}        & 32 & Kanal-Adresse Bit 0 \\
			\texttt{A1}        & 4  & Kanal-Adresse Bit 1 \\ \bottomrule
		\end{tabular}
	\end{table}

	\subsubsection{Serielles Leseprotokoll}

	Der ADS1234 signalisiert eine abgeschlossene Wandlung, indem er die \texttt{DRDY}-Leitung auf \texttt{LOW} zieht.
	Die Software wartet aktiv auf diesen Zustand und liest anschliessend 24 Bit MSB-first aus, indem für jedes Bit ein \texttt{SCLK}-Impuls erzeugt und der Pegel von \texttt{DOUT} abgetastet wird.
	Nach dem 24. Bit wird gemäss Datenblatt (Figure~7-10) ein 25.~Taktimpuls gesendet, der \texttt{DRDY} wieder auf \texttt{HIGH} zwingt und damit den ADC für die nächste Wandlung freigibt.
	Da der ADS1234 ein Zweierkomplement-Format mit 24~Bit Breite verwendet, wird der Rohwert abschliessend vorzeichenrichtig auf 32~Bit erweitert:

	\begin{verbatim}
	if (raw & 0x800000) raw |= 0xFF000000;
	return (int32_t)raw;
	\end{verbatim}

	\subsubsection{Kanalwechsel CH4 $\rightarrow$ CH1: PDWN-Workaround}

	Der ADS1234 kodiert den aktiven Kanal über die beiden Adressleitungen als 2-Bit-Binärzahl: Kanal~1 entspricht \texttt{A0=0, A1=0}, Kanal~4 entspricht \texttt{A0=1, A1=1}.
	Der Chip erkennt einen Kanalwechsel intern anhand steigender Flanken auf \texttt{A0} oder \texttt{A1}.
	Beim Wechsel von Kanal~4 nach Kanal~1 fallen jedoch \emph{beide} Adressleitungen von \texttt{HIGH} auf \texttt{LOW} -- es entsteht keine einzige steigende Flanke.
	Der ADS1234 erkennt diesen Wechsel daher nicht als Kanalumschaltung und setzt \texttt{DRDY} nicht zurück.
	Die Leitung bleibt \texttt{LOW}, was der Software fälschlicherweise signalisiert, dass bereits ein gültiger Messwert des neuen Kanals vorliegt, obwohl noch der alte Kanal aktiv ist.

	Als Lösung wird nach dem Setzen der Adressleitungen auf Kanal~1 geprüft, ob \texttt{DRDY} bereits \texttt{LOW} ist.
	Trifft dies zu, wird ein \texttt{PDWN}-Puls von \SI{500}{\micro\second} ausgelöst, der den ADC in den Power-Down-Zustand versetzt und anschliessend neu startet.
	Dieser Reset zwingt den Chip, den ausgewählten Kanal neu zu initialisieren und eine frische Wandlung durchzuführen, bevor \texttt{DRDY} das nächste Mal \texttt{LOW} wird.

	\begin{verbatim}
	int32_t readChannel(int ch) {
	    setChannel(ch);
	    if (ch == 1 && digitalRead(PIN_DRDY_DOUT) == LOW) {
	        digitalWrite(PIN_PDWN, LOW);
	        delayMicroseconds(500);
	        digitalWrite(PIN_PDWN, HIGH);
	    }
	    return readADS1234();
	}
	\end{verbatim}

	\subsubsection{Kalibrierung}

	Die Kalibrierung läuft in drei aufeinanderfolgenden Schritten ab und wird einmalig beim Start der Firmware über die serielle Schnittstelle ausgelöst.
	Alle Messwerte sind Mittelwerte aus \texttt{CAL\_SAMPLES}~=~5 Wandlungen pro Kanal, was bei 10~SPS einer Messzeit von rund 0,4~Sekunden pro Kanal entspricht.

	\paragraph{Schritt 1 -- Nullpunkt (Zero Offset)}
	Die Waage wird vollständig entlastet und für jeden der vier Kanäle ein gemittelter Rohwert erfasst, der als Nullpunkt \texttt{zeroOffset[ch]} gespeichert wird.
	Dieser Wert kompensiert thermische und mechanische Vorlast der Wägezellen sowie Offset-Fehler des ADC.

	\paragraph{Schritt 2 -- Empfindlichkeit (Span)}
	Ein bekanntes Kalibriergewicht von \SI{3000}{\gram} wird mittig auf die Waagplattform gelegt.
	Für jeden Kanal wird die Abweichung vom Nullpunkt berechnet:
	\begin{equation}
		\Delta_i = \overline{x}_i - \text{zeroOffset}_i
	\end{equation}
	Die vier Wägezellen sind in X-Formation an den Ecken der Plattform angebracht.
	Bei dieser Anordnung werden die zwei diagonal gegenüberliegenden Zellen beim Belasten auf Zug, die anderen zwei auf Druck beansprucht, was zu einem mechanisch bedingten Vorzeichenwechsel im ADC-Signal führt.
	Das Vorzeichen von $\Delta_i$ wird daher pro Kanal automatisch erfasst und als \texttt{channelSign} gespeichert.
	Der initiale Skalierungsfaktor ergibt sich aus der Gesamtauslenkung über alle vier Kanäle:
	\begin{equation}
		s_{\text{init}} = \frac{m_{\text{cal}}}{\sum_{i=1}^{4} |\Delta_i|}
		\qquad \Rightarrow \qquad
		\text{spanFactor}_i = \text{channelSign}_i \cdot s_{\text{init}}
	\end{equation}
	Damit trägt jeder Kanal anteilig zur Gesamtmasse bei, wobei die Polarität korrekt berücksichtigt wird.

	\paragraph{Schritt 3 -- Ecklastkompensation (Corner Trim)}
	Bei einer Plattformwaage mit vier Wägezellen in den Ecken führen fertigungsbedingte Toleranzen der Sensoren dazu, dass ein exzentrisch aufgelegtes Gewicht -- je nach Position -- unterschiedliche Gesamtmassen anzeigt.
	Um dies zu kompensieren, wird das Kalibriergewicht nacheinander über jede der vier Ecken gestellt.

	Bei jeder Eckenbelastung ermittelt die Software denjenigen Kanal mit der grössten absoluten Auslenkung (\emph{dominanter Kanal}), da dieser am stärksten von der Lastposition beeinflusst wird.
	Die aktuelle Gesamtanzeige $m_{\text{meas}}$ ergibt sich aus den Beiträgen aller vier Kanäle:
	\begin{equation}
		m_{\text{meas}} = \sum_{i=1}^{4} \text{spanFactor}_i \cdot \Delta_i
	\end{equation}
	Der Skalierungsfaktor des dominanten Kanals~$d$ wird anschliessend so angepasst, dass die Anzeige exakt dem Kalibriergewicht entspricht, während die Faktoren der übrigen Kanäle unverändert bleiben:
	\begin{equation}
		\text{spanFactor}_d = \frac{m_{\text{cal}} - \sum_{i \neq d} \text{spanFactor}_i \cdot \Delta_i}{\Delta_d}
	\end{equation}
	Ein einziger Durchlauf über alle vier Ecken ist für typische DMS-Toleranzen von unter 5\,\% ausreichend.

	\paragraph{Gewichtsberechnung im Betrieb}
	Nach abgeschlossener Kalibrierung berechnet sich die angezeigte Masse in jedem Messtakt als:
	\begin{equation}
		m = \sum_{i=1}^{4} \text{spanFactor}_i \cdot \bigl(\text{raw}_i - \text{zeroOffset}_i\bigr)
	\end{equation}

	% ============================================================
	\section{Tests und Validierung}
	% ============================================================

	\subsection{Strommessungen und Low-Power-Optimierung}

	Ein zentrales Ziel dieser Arbeit ist die Minimierung des Stromverbrauchs, um eine möglichst lange Batterielaufzeit zu erreichen.
	Die folgenden Messungen dokumentieren den Ladungsverbrauch in den einzelnen Betriebsphasen und quantifizieren den Effekt der implementierten Optimierungen.
	Tabelle~\ref{tab:strom-uebersicht} fasst die verbrauchte Ladung pro Phase zusammen -- jeweils vor und nach der Optimierung.

	
	\subsubsection{Setup}
	Die Strommessungen wurden mit dem \textit{Nordic Semiconductor Power Profiler Kit~II} (PPK2) in Verbindung mit der Software \textit{nRF Connect for Desktop} (Power Profiler App) durchgeführt.
	Das PPK2 wurde als Amperemeter in den Versorgungspfad des Zielsystems eingeschleift.
	Das Gerät erfasst den Strom mit einer Abtastrate von bis zu $100{,}000\,\text{Samples/s}$, was eine präzise zeitliche Auflösung auch kurzer Stromspitzen ermöglicht.

	In der nRF Connect Power Profiler App lassen sich Zeitfenster (\emph{Selektionen}) manuell markieren, woraufhin die Software automatisch den mittleren Strom, die Dauer sowie den gesamten Ladungsverbrauch der Selektion in Millicoulomb (\si{mC}) und Milliampere-Stunden (\si{mAh}) berechnet.
	Alle in dieser Arbeit angegebenen Ladungswerte basieren auf solchen Selektionen aus aufgezeichneten Stromprofilen.
	
	\subsubsection{Ausgangslage}
	Abbildung~\ref{fig:stromverbrauch-ausgangslage} zeigt den Stromverlauf eines vollständigen Messzyklus in der Ausgangslage, wie er direkt aus der Vorgängerarbeit übernommen wurde.
	In dieser Vorarbeit stand die funktionale Korrektheit im Vordergrund -- eine gezielte Optimierung des Stromverbrauchs hatte keine Priorität.

	Im Profil sind drei charakteristische Abschnitte erkennbar: ein kurzer initialer Peak während Boot und Initialisierung, ein ruhigeres Plateau während der ADC-Messung sowie ein erhöhter und zeitlich ausgedehnter Abschnitt während der BLE-Kommunikation.
	Die gesamte Aktivphase erstreckt sich über ca.\ \textbf{10\,s} mit einem Gesamtladungsverbrauch von \textbf{0.314\,mAh}, bevor das System in den Deep Sleep übergeht.
	
	
	\begin{figure}[h]
	\centering
	\includegraphics[width=0.85\linewidth]{"Images/Stromverbrauch/Stromverbrauch vorher/Stromverbrauch Ausgangslage"}
	\caption{Stromverbrauch -- Ausgangslage, ein vollständiger Messzyklus}
	\label{fig:stromverbrauch-ausgangslage}
	\end{figure}

	\subsubsection{Übersicht des Zyklus der Waage}
	\begin{enumerate}
	\item \textbf{Boot (0\,s):} Der ESP32 startet nach dem Deep Sleep neu. Die State Machine wird initialisiert.
	\item \textbf{ADC-Initialisierung (0\,s):} Der ADS1234 wird eingeschaltet und auf Betriebsbereitschaft geprüft (\texttt{ADC ready}).
	\item \textbf{Messung (1--2\,s):} Das Gewicht wird erfasst und auf Änderung geprüft (\texttt{Weight changed: true}).
	\item \textbf{BLE-Initialisierung (2--3\,s):} Der BLE-Stack wird gestartet und Advertisement begonnen.
	\item \textbf{Datenübertragung (3--10\,s):} Das System wartet auf die Verbindung des Gateways (~6\,s Wartezeit), überträgt den Messwert per Indication und beendet die Verbindung.
	\item \textbf{Deep Sleep (ab 10\,s):} Das System tritt in den Deep Sleep ein und bleibt bis zum nächsten Messzyklus inaktiv.
	\end{enumerate}

	
		\begin{table}[h]
		\centering
		\caption{Ladungsverbrauch pro Betriebsphase: Ausgangslage vs.\ Optimiert}
		\label{tab:strom-uebersicht}
		\begin{tabular}{@{}lccc@{}}
			\toprule
			Betriebsphase & Ausgangslage [mAh] & Optimiert [mAh] & Einsparung [mAh] \\ \midrule
			Deep Sleep                              & --    & --    & -- \\
			Ramp Up (Boot, Init, ADC)               & 0.123 & --    & -- \\
			BLE-Init \& Advertisement$^{*}$         & --    & --    & -- \\
			BLE-Verbindung \& Datenübertragung$^{*}$& 0.194 & --    & -- \\
			\midrule
			Aktivphase gesamt                       & 0.314 & --    & -- \\ \bottomrule
		\end{tabular}
		\\[4pt]
		\small $^{*}$In der Ausgangslage gemeinsam als eine Selektion gemessen (0.194\,mAh total).
	\end{table}

	Die gesamte Aktivphase beträgt in der Ausgangslage ca.\ \textbf{10 Sekunden} pro Zyklus mit einem Ladungsverbrauch von \textbf{0.314\,mAh} (Abbildung~\ref{fig:stromverbrauch-ausgangslage}).

	Die nachfolgenden Abschnitte beschreiben jede Betriebsphase im Detail -- mit Stromprofilen der Ausgangslage sowie der optimierten Version -- und erläutern, welche Massnahmen zur jeweiligen Einsparung geführt haben:

	\begin{enumerate}
		\item \textbf{Deep Sleep} (Abschnitt~\ref{sec:deep-sleep}): \\
		Im Deep Sleep dominiert der Ruhestrom des Systems.
		Durch gezieltes Abschalten externer Peripherie (ADS1234, Wägezellen-Versorgung) vor dem Schlafeingang sowie die Wahl eines LDO mit minimalem Quiescent Current wurde der Sockellader auf ein Minimum reduziert.

		\item \textbf{Ramp Up -- Boot, Initialisierung und ADC} (Abschnitt~\ref{sec:ramp-up}): \\
		Diese Phase umfasst den ESP32-Bootvorgang, die State-Machine-Initialisierung sowie die ADS1234-Messung.
		Die Ausgangslage verbrauchte \textbf{0.123\,mAh}.
		Optimierungspotenzial liegt unter anderem in der Reduktion der Boot-Zeit sowie in der Minimierung der ADC-Einschaltdauer.

		\item \textbf{ADC} (Abschnitt~\ref{sec:ramp-up}): \\
			In der Ausgangslage verwendete die Vorgängerarbeit den \textbf{HX711} -- einen einkanaligen 24-Bit-ADC, an den alle vier Wägezellen als gemeinsame Brücke angeschlossen waren.
			Die Messung erfasste 20 Samples auf diesem einzigen Kanal; zwischen jedem Sample wurde ein blockierendes \texttt{delay(50\,ms)} ausgeführt, was den ESP32-Kern für die gesamte Wartezeit vollaktiv liess.

			\begin{figure}[h]
			\centering
			\includegraphics[width=0.85\linewidth]{"Images/Stromverbrauch/Stromverbrauch vorher/Strom ADC"}
			\caption{Stromverlauf -- ADC (Ausgangslage)}
			\label{fig:strom-bootup}
		\end{figure}

			Im optimierten Design werden die vier Kanäle des ADS1234 sequenziell ausgelesen; pro Kanal werden \textbf{5 Samples} erfasst.
			Nach jedem abgeschlossenen Sample kehrt der ESP32 sofort in den Light Sleep zurück und wird durch einen GPIO-Interrupt auf dem \texttt{DRDY}-Pin geweckt, sobald der ADS1234 die nächste Wandlung bereitstellt.
			Im Profil sind diese fünf kurzen Aktivierungsspitzen pro Kanalgruppe deutlich erkennbar.

			Zwischen den vier Kanalgruppen zeigt sich jeweils eine Pause von ca.\ \textbf{400\,ms}: Dies ist die physikalisch bedingte Einschwingzeit des internen Sigma-Delta-Filters nach einem Kanalwechsel.
			Der ADS1234 integriert einen digitalen sinc\textsuperscript{3}-Filter; dessen Impulsantwort erstreckt sich über 3--4 Ausgangsperioden, was bei der eingestellten Datenrate von 10\,Hz rund 300--400\,ms entspricht.
			Erst nach dieser Zeit sind die Ausgabewerte frei von Einflüssen des vorangegangenen Kanals.
			Auch diese Wartezeit verbringt der ESP32 im Light Sleep -- der Kern schläft, bis \texttt{DRDY} das Ende der Einschwingphase signalisiert.

			Der Gesamtladungsverbrauch der ADC-Phase sinkt von ca.\ \textbf{190\,mC} auf ca.\ \textbf{100\,mC} -- eine Einsparung von rund \textbf{47\,\%}.

			\begin{figure}[h]
			\centering
			\includegraphics[width=0.85\linewidth]{"Images/Stromverbrauch/Nachher/ADC"}
			\caption{Stromverlauf -- ADC (Optimiert)}
			\label{fig:strom-bootup}
		\end{figure}
		
		

		\item \textbf{BLE-Initialisierung und Advertisement} (Abschnitt~\ref{sec:ble-init}): \\
		Der NimBLE-Stack wird gestartet und das Advertisement-Paket ausgesendet.
		Durch den Wechsel von der Standard-Bluedroid-Implementierung auf den schlanken NimBLE-Stack konnte der Speicherverbrauch und die Initialisierungszeit signifikant reduziert werden.
		
			\begin{figure}[h]
			\centering
			\includegraphics[width=0.85\linewidth]{"Images/Stromverbrauch/Stromverbrauch vorher/subscription"}
			\caption{Stromverlauf -- BLE Initialisierung und Subscription (Ausgangslage)}
			\label{fig:strom-bootup}
		\end{figure}
		
			\begin{figure}[h]
			\centering
			\includegraphics[width=0.85\linewidth]{"Images/Stromverbrauch/Nachher/BLE_GANZ"}
			\caption{Stromverlauf -- BLE Initialisierung und Advertisement (Optimiert)}
			\label{fig:strom-bootup}
		\end{figure}

		\item \textbf{BLE-Verbindung und Datenübertragung} (Abschnitt~\ref{sec:ble-data}): \\
		Das Gateway baut die Verbindung auf, der Messwert wird per Indication übertragen und die Verbindung anschliessend getrennt.
		In der Ausgangslage betrug der Ladungsverbrauch dieser Phase (gemeinsam mit BLE-Init gemessen) \textbf{0.194\,mAh}.
		Durch Optimierung der Connection-Parameter und Reduktion der Wartezeit auf den Gateway konnte dieser Wert gesenkt werden.
	\end{enumerate}

	\subsubsection{Deep Sleep}
	\label{sec:deep-sleep}

	Im Deep Sleep ist der ESP32 vollständig inaktiv; lediglich der RTC-Timer läuft weiter, um das System nach dem konfigurierten Intervall aufzuwecken.
	Der Ruhestrom setzt sich in dieser Phase aus dem Quiescent Current des LDO (RT9080: $2\,\mu\text{A}$), dem Deep-Sleep-Strom des ESP32-C3 sowie allfälligen Leckströmen durch externe Peripherie zusammen.

	% TODO: Screenshot Deep Sleep Ausgangslage einfügen
	% \begin{figure}[h]
	%     \centering
	%     \includegraphics[width=0.85\linewidth]{"Images/Stromverbrauch/Stromverbrauch vorher/DeepSleep"}
	%     \caption{Stromverlauf -- Deep Sleep (Ausgangslage)}
	%     \label{fig:strom-deepsleep-vorher}
	% \end{figure}

	In der Ausgangslage blieben der ADS1234 und die Wägezellen-Versorgungsleitung (\texttt{DMS\_PWR}) beim Eintreten in den Deep Sleep aktiv, was zu einem erhöhten Sockellader führte.
	Durch das explizite Abschalten dieser Peripherie vor dem \texttt{esp\_deep\_sleep\_start()}-Aufruf konnte der Ruhestrom auf den spezifizierten Minimalwert abgesenkt werden.

	% TODO: Screenshot Deep Sleep Optimiert einfügen
	% \begin{figure}[h]
	%     \centering
	%     \includegraphics[width=0.85\linewidth]{"Images/Stromverbrauch/Stromverbrauch nachher/DeepSleep"}
	%     \caption{Stromverlauf -- Deep Sleep (Optimiert)}
	%     \label{fig:strom-deepsleep-nachher}
	% \end{figure}

	\subsubsection{Ramp Up: Boot, Initialisierung und ADC}
	\label{sec:ramp-up}

	Nach dem Aufwachen aus dem Deep Sleep durchläuft das System zuerst die Boot- und Initialisierungsphase.
	Abbildung~\ref{fig:strom-bootup} zeigt den Stromverlauf dieser Phase in der Ausgangslage.
	Die Selektion umfasst \textbf{2.581\,s} mit einem Ladungsverbrauch von \textbf{253.78\,mC} (entspricht \textbf{0.070\,mAh}).

	\begin{figure}[h]
		\centering
		\includegraphics[width=0.85\linewidth]{"Images/Stromverbrauch/Stromverbrauch vorher/bootup"}
		\caption{Stromverlauf -- Boot \& Initialisierung (Ausgangslage)}
		\label{fig:strom-bootup}
	\end{figure}

	Im Anschluss erfolgt die ADC-Messung.
	Abbildung~\ref{fig:strom-adc} zeigt den Verlauf während des ADC-Samplings.
	Die Selektion umfasst \textbf{1.769\,s} mit einem Ladungsverbrauch von \textbf{189.97\,mC} (entspricht \textbf{0.053\,mAh}).

	\begin{figure}[h]
		\centering
		\includegraphics[width=0.85\linewidth]{"Images/Stromverbrauch/Stromverbrauch vorher/Strom ADC"}
		\caption{Stromverlauf -- ADC-Messung (Ausgangslage)}
		\label{fig:strom-adc}
	\end{figure}

	% TODO: Screenshot Ramp Up Optimiert einfügen
	% \begin{figure}[h]
	%     \centering
	%     \includegraphics[width=0.85\linewidth]{"Images/Stromverbrauch/Stromverbrauch nachher/bootup"}
	%     \caption{Stromverlauf -- Boot \& Initialisierung (Optimiert)}
	%     \label{fig:strom-bootup-nachher}
	% \end{figure}

	% TODO: Screenshot ADC Optimiert einfügen
	% \begin{figure}[h]
	%     \centering
	%     \includegraphics[width=0.85\linewidth]{"Images/Stromverbrauch/Stromverbrauch nachher/Strom ADC"}
	%     \caption{Stromverlauf -- ADC-Messung (Optimiert)}
	%     \label{fig:strom-adc-nachher}
	% \end{figure}

	\subsubsection{BLE-Initialisierung und Advertisement}
	\label{sec:ble-init}
	% Hier Text einfügen (Screenshots BLE Init und Advertisement -- Ausgangslage und Optimiert)

	\subsubsection{BLE-Verbindung und Datenübertragung}
	\label{sec:ble-data}

	In der Ausgangslage wurden BLE-Init, Advertisement, Verbindungsaufbau und Datenübertragung als zusammenhängende Phase gemessen.
	Abbildung~\ref{fig:strom-ble} zeigt die gesamte BLE-Phase.
	Die Selektion umfasst \textbf{5.642\,s} mit einem Ladungsverbrauch von \textbf{700\,mC} (entspricht \textbf{0.194\,mAh}).

	\begin{figure}[h]
		\centering
		\includegraphics[width=0.85\linewidth]{"Images/Stromverbrauch/Stromverbrauch vorher/subscription"}
		\caption{Stromverlauf -- BLE-Phase gesamt (Ausgangslage)}
		\label{fig:strom-ble}
	\end{figure}

	% Hier Text einfügen (optimierte Messungen BLE-Verbindung -- Screenshot nachher)

	% ============================================================
	\section{Fazit, Diskussion und Ausblick}
	% ============================================================
	% Hier Text einfügen

	\bibliographystyle{plain}
	\bibliography{quellen}

\end{document}
